\numberwithin{equation}{section}

\newtheorem{Theorem}{Theorem}[section]
\newtheorem{Corollary}[Theorem]{Corollary}
\newtheorem{Lemma}[Theorem]{Lemma}
\newtheorem{Proposition}[Theorem]{Proposition}
 { \theoremstyle{definition}
\newtheorem{Definition}[Theorem]{Definition}
\newtheorem{Note}[Theorem]{Note}
\newtheorem{Example}[Theorem]{Example}
\newtheorem{Remark}[Theorem]{Remark} }


\DeclareMathOperator{\Id}{Id}
\DeclareMathOperator{\Span}{Span}
\DeclareMathOperator{\Ann}{Ann}
\DeclareMathOperator{\End}{End}
\DeclareMathOperator{\Hom}{Hom}
\DeclareMathOperator{\Aut}{Aut}

\usepackage{enumitem}
\usepackage{mathrsfs}
% \iffalse meta-comment
%
% File: romanbar.dtx
% Version: 2025-01-28 v1.0g
% Info: Roman numbers with bars
%
% Copyright (C) 2011 - 2025 by
%    H.-Martin Münch <Martin dot Muench at Uni-Bonn dot de>
%
% The original code for the bars was written by Prof. Enrico Gregorio
% (https://profs.sci.univr.it/~gregorio/), and published at
% https://tex.stackexchange.com/a/24084
% I (H.-Martin Münch) submitted this package to CTAN
% (after diverse updates of the code and writing a documentation).
%
% This work may be distributed and/or modified under the
% conditions of the LaTeX Project Public License, either
% version 1.3c of this license or (at your option) any later
% version. This version of this license is in
%    https://www.latex-project.org/lppl/lppl-1-3c.txt
% and the latest version of this license is in
%    https://www.latex-project.org/lppl.txt
% and version 1.3c or later is part of all distributions of
% LaTeX version 2005-12-01 or later.
%
% This work has the LPPL maintenance status "maintained".
%
% The Current Maintainer of this work is H.-Martin Münch.
%
% This work consists of the main source file romanbar.dtx,
% the README, and the derived files
% CTAN romanbar1.0g (H.-Martin Muench), LPPL, embedded macro code.
\makeatletter

%    \end{macrocode}
%
%    \begin{macrocode}
%    \end{macrocode}
%
% A short description of the \xpackage{romanbar} package:
%
%    \begin{macrocode}
%% Allows to write Roman numbers (or any other text) with bars
%% and to convert Arabic numbers into Roman ones.
%    \end{macrocode}
%
% There are no options to be processed.
%
%    \begin{macrocode}

\RequirePackage{trimclip}[2020/08/19]% v1.2

\newcounter{romanbarcount}
\newdimen\romanbardimH
\newdimen\romanbardimD
\newdimen\romanbardimW

%    \end{macrocode}
%    \begin{macro}{\gobbleminus}
%
% A command to gobble a possible minus sign is needed and
% |\gobbleminus| from \url{https://texfaq.org/FAQ-isitanum}
% is used.
%
%    \begin{macrocode}
\def\gobbleminus#1{\ifx-#1\else#1\fi}
%% from https://texfaq.org/FAQ-isitanum

%    \end{macrocode}
%    \end{macro}
%    \begin{macro}{\ifnumeric}
%
% We want to test whether the argument passed to |\Romanbar| is numeric
% and therefore define |\ifnumeric{test}{true}{false}|,
% where |test| is to be analysed, |true| is the code to be executed,
% if |test| is numeric, and |false| is the code to be executed,
% if |test| is not numeric. This is done similar to
% \url{https://tex.stackexchange.com/a/17119}.
%
%    \begin{macrocode}
\newcommand{\ifnumeric}[3]{%
%% similar to https://tex.stackexchange.com/a/17119
\sbox\z@{\c@romanbarcount=0\gobbleminus#1\relax}%
\ifdim\wd0>\z@\relax#3% is not numeric
\else#2% is numeric
\fi%
}

%    \end{macrocode}
%    \end{macro}
%    \begin{macro}{\r@iseL}
%
% We define |\r@iseL| to raise any L~(50), otherwise a good deal
% of the horizontal part of the letter would be swallowed by the lower bar
% under the \textquotedblleft number\textquotedblright .
%
%    \begin{macrocode}
\def\r@iseL#1{\ifx @#1% then terminate
\else%
\if L#1\raisebox{0.05ex}{\begin{clipbox}{0pt 0pt 0pt 0.05ex}L\end{clipbox}}%
\else #1\fi%
\expandafter\r@iseL%
\fi%
}

%    \end{macrocode}
%    \end{macro}
%    \begin{macro}{\Romanbar}
%
% We define the |\Romanbar| command.
%
%    \begin{macrocode}
\DeclareRobustCommand{\Romanbar}[1]{%
\ifnumeric{#1}{% is numeric
%    \end{macrocode}
%
% If its parameter is numeric, the temporary counter |romanbarcount| is set to the value.
%
%    \begin{macrocode}
\setcounter{romanbarcount}{#1}%
%    \end{macrocode}
%
% Regularly there are no non-positive Roman numbers, therefore we do some trick
% similar to |\XRoman| from the \url{https://ctan.org/pkg/pageslts} package.
%
%    \begin{macrocode}
%% similar to \XRoman from the https://ctan.org/pkg/pageslts package
\ifnum\theromanbarcount<1%
  \ifnum #1>0\relax%
    \Roman@bar{\Roman{#1}}%
  \else%
    \ifnum #1<0\relax%
      -\Roman@bar{\expandafter\@slowromancap\romannumeral\number-#1@}%
    \else%
      \Roman@bar{0}%
    \fi%
  \fi%
\else\Roman@bar{\@Roman\c@romanbarcount}%
\fi%
}{% is not numeric
\Roman@bar{#1}%
}%
}

%    \end{macrocode}
%    \end{macro}
%    \begin{macro}{\Roman@bar}
%
% Whether it is numeric or not, in the end |\Roman@bar| is applied to it.
%
%    \begin{macrocode}
\DeclareRobustCommand{\Roman@bar}[1]{% #1 is in Roman, i.e. MMXXIII
%% (or any other text, "Caesar" would work, too);
%% similar to code of  by Prof. Enrico Gregorio (egreg) at
%% https://tex.stackexchange.com/a/24084
\@bsphack%
\edef\romanbartmp{#1}%
%% height:
\sbox0{\textsf{\romanbartmp}}%
%% top line:
%% + top of top line:
\setlength{\romanbardimH}{\dimexpr\ht0+0.05ex\relax}%
%% + bottom of top line:
\setlength{\romanbardimD}{\dimexpr-\ht0+0.05ex\relax}%
%% width:
\sbox0{\textsf{\expandafter\r@iseL\romanbartmp @}}%
\setlength{\romanbardimW}{\dimexpr\wd0+1pt\relax}%
\@esphack%
\makebox[\romanbardimW]{%
\rlap{\vrule width\romanbardimW height\romanbardimH depth\romanbardimD}% top line
\rlap{\vrule width\romanbardimW height 0.06ex       depth 0.03ex}%    bottom line
\kern0.5pt\textsf{\expandafter\r@iseL\romanbartmp @}\kern0.5pt}%
}

%    \end{macrocode}
%    \end{macro}
%    \begin{macro}{\romannum}
%
% While it is not needed for the bars, it is sometimes asked for a command
% to convert an Arabic number into a Roman one. |\roman{...}| can only be applied
% to a counter, i.\,e.~|\roman{25}| does not work, but
%\begin{verbatim}
%\newcounter{examplecounter}
%\setcounter{examplecounter}{25}
%\roman{examplecounter}
%\end{verbatim}
% would be necessary. Further non-negative values would not work.
% |\romannum| accepts any number as argument. (If the argument is no number,
% an error message is given and the argument printed (executed)
% without applying anything to it.)
%
%    \begin{macrocode}
\DeclareRobustCommand{\romannum}[1]{%
\ifnumeric{#1}{% is numeric
\setcounter{romanbarcount}{#1}%
%% similar to \xroman from the https://ctan.org/pkg/pageslts package
\ifnum\c@romanbarcount<1\relax%
  \ifnum #1>0\relax%
    \roman{#1}%
  \else%
    \ifnum #1<0\relax%
      -\romannumeral\number-#1@%
    \else%
      0%
    \fi%
  \fi%
\else\@roman\c@romanbarcount%
\fi%
}{% is not numeric
\PackageError{romanbar}{%
Argument of \string\romannum\space is not a number}{%
The command \string\romannum\space converts an Arabic number into a %
lower-case Roman one,\MessageBreak%
but the used argument of \string\romannum\space is not an Arabic number %
but\MessageBreak%
`#1`,\MessageBreak%
which will now be printed unchanged.}%
#1%
}%
}

%    \end{macrocode}
%    \end{macro}
%    \begin{macro}{\Romannum}
%
% The same for upper-case Roman numbers, |\Romannum| instead of |\Roman|:
%
%    \begin{macrocode}
\DeclareRobustCommand{\Romannum}[1]{%
\ifnumeric{#1}{% is numeric
\setcounter{romanbarcount}{#1}%
%% similar to \XRoman from the http://ctan.org/pkg/pageslts package
\ifnum\c@romanbarcount<1\relax%
  \ifnum #1>0\relax%
    \Roman{#1}%
  \else%
    \ifnum #1<0\relax%
      -\expandafter\@slowromancap\romannumeral\number-#1@%
    \else%
      0%
    \fi%
  \fi%
\else\@Roman\c@romanbarcount%
\fi%
}{% is not numeric
\PackageError{romanbar}{%
Argument of \string\Romannum\space is not a number}{%
The command \string\Romannum\space converts an Arabic number into an %
upper-case Roman one,\MessageBreak%
but the used argument of \string\Romannum\space is not an Arabic number %
but\MessageBreak%
`#1`.}%
#1%
}%
}

%    \end{macrocode}
%    \end{macro}
%    \begin{macrocode}
\makeatother
\newcommand{\II}{\mathrm{I\hspace{-.1em}I}}

\allowdisplaybreaks[4]
\pdfstringdefDisableCommands{\def\boldsymbol#1{#1}}
